\chapter{Sistemas inversos, representación de Weyl--Heisenberg y cociente predictivo}
\label{chap:numero-heisenberg-d108}
\index{número HMT}
\index{Heisenberg!grupo finito}

La construcción distingue tres niveles: el catálogo finito de semillas, el
espacio de estados que conserva el acarreo y la orientación, y el límite
inverso de las prolongaciones compatibles. Esta estratificación permite
formalizar la contracción de los intervalos cilíndricos hasta un punto real,
la persistencia de datos de trayectoria en una misma fibra de evaluación y
la representación finita de Weyl--Heisenberg asociada a la forma alternante
del toro nonádico que soporta los observables aditivo y multiplicativo.

Los símbolos \(V_t\), \(X_t\) y \(\rho_{t+1,t}\) empleados aquí son los
mismos conjuntos de estados locales, prefijos admisibles y truncamientos
introducidos en el \cref{chap:numero-medicion}, con el cambio puramente
notacional \(t=m\) y \(\rho_{n,m}=p_{n,m}\). Esta sección explicita sus
coordenadas y no introduce un segundo sistema inverso.

\section{Proyección observable y espacio completo de estados del TPK}

La proyección observable del sistema de transporte de fase (TPK) es la
aplicación finita de codificación
\[
104\,976\longrightarrow468\longrightarrow243.
\]
que clasifica semillas, firmas y palabras \(w_6\). El espacio completo de
estados del TPK conserva además
\[
\begin{gathered}
\text{orientación bidireccional},\quad
\text{estado Hensel},\quad
\text{cilindro y acarreo},\\
\text{firma conjunta de frontera},\quad
\text{fase nonádica visible }g_{\rm vis}(m).
\end{gathered}
\]
Dos candidatos con una misma firma local pueden pertenecer a clases distintas
cuando se conserva la frontera siguiente. No introducimos una segunda tupla
de estado. Para
\[
x_m=(q_m,o_m,h_m,c_m,\sigma_m,C_m,b_m,\lambda_m)\in\XTPK^{(m)}
\]
usamos la carta de coordenadas
\[
\chi_m(x_m)=
\bigl(
B_m,h_m,(C_m,b_m),\sigma_m,
o_m,\mu(q_m),g_{\rm vis}(m),\lambda_m
\bigr),
\]
con
\[
B_m=(u_m^\pi,u_m^e,u_m^\varphi)\in(\mathbb F_3^6)^3.
\]
Aquí \(h_m\) es la extensión de Hensel; \((C_m,b_m)\), el cilindro y su
frontera; \(o_m,\mu(q_m)\), la orientación y el observable bidireccional; y
\[
g_{\rm vis}(m)=1+((m-1)\bmod9)\in\{1,\ldots,9\},
\]
la etiqueta visible de fase nonádica. La coordenada de clase correspondiente
es \(g_0(\theta_m)\in\ZZ/9\ZZ\), con \(\theta_m=[m-1]_{108}\). La firma conjunta
es la coordenada ya tipada por el núcleo de fase,
\[
\sigma_m=\operatorname{Sig}_{q_m}(h_m,c_m,\lambda_m)
=(q_\partial,a_\partial,c_\partial,
\operatorname{colw},\rho_W,r_\partial).
\]
La firma conjunta no se reduce a tres firmas de canal independientes, pues la
selección asociada a la involución \(\pi/e\) y al eje \(\varphi\) depende de
datos correlacionados. Dos nodos con el mismo \(B_m\) representan estados
distintos cuando difieren en extensión, cilindro, acarreo, firma, orientación
o registro de transporte.

\subsection{Compatibilidad exacta de los cilindros bajo truncación y prolongación nonádica}

Sea \(N\) el entero de un prefijo ternario de longitud \(L\), y sea \(D\)
el entero obtenido al concatenar sus \(b\) tríadas decimales. Definimos
\[
A=N1000^b-D3^L,
\qquad
B=(N+1)1000^b-D3^L.
\]
Los cilindros se intersecan si y sólo si
\[
\boxed{A<3^L,\qquad B>0.}
\]
Al añadir un bloque ternario \(u\in\{0,\ldots,728\}\), se tiene
\[
N'=729N+u.
\]
Si no aparece una nueva tríada decimal,
\[
A'=729A+u1000^b.
\]
Si aparece \(\delta\), de modo que \(D'=1000D+\delta\), entonces
\[
A'=729000A+u1000^{b+1}-\delta\,729\,3^L.
\]
Si \(\Delta b\in\{0,1\}\) registra la aparición de esa nueva tríada, en
ambos casos
\[
B'=A'+1000^{b+\Delta b}.
\]
La transición, fijado el par admisible \((u,\delta)\), no consulta un valor
objetivo. La filtración conjunta por intersección de cilindros y retorno
nonádico actúa en un nivel distinto: selecciona los pares de frontera que
admiten prolongación compatible.

\subsection{Registro cíclico de 108 etapas y su proyección reducida}

Denominamos \emph{registro cíclico de 108 etapas} a la sucesión de
\(108\) pasos agrupados en doce sectores de nueve pasos, con la convención de
índice fijada en el \cref{chap:tpk-definicion}.

Sea \(\XTPK^{[108]}\subseteq\XTPK\) el dominio de estados completos cuyo
registro \(\lambda\) contiene los \(108\) pasos declarados. Denotemos por
\(\mathcal C_{108}^{\rm red}\) el conjunto de registros que resultan al
conservar sólo los campos reducidos, y por \(\mathcal U_{12}^{\rm sgn}\) el
conjunto de dodecafases firmadas admisibles. La instancia distingue dos mapas
declarados:
\[
F:\XTPK^{[108]}\longrightarrow\mathcal C_{108}^{\rm red},
\qquad
R:\XTPK^{[108]}\longrightarrow\mathcal U_{12}^{\rm sgn}.
\]
El mapa \(F\) retiene sólo los campos del registro reducido, mientras que
\(R\) evalúa la dodecafase firmada. La existencia de una aplicación \(\Xi\)
tal que \(R=\Xi\circ F\) equivale a la factorización de \(R\) a través de una
proyección que puede eliminar la coordenada del observable, la orientación,
el acarreo o la frontera. Esta factorización constituye una proposición
adicional y no forma parte implícita de la reconstrucción sobre el espacio
completo de estados.

El registro reducido es, por definición, una proyección del estado completo,
no el estado de partida. Los resultados posteriores usan \(R\) sobre el
estado completo y no presuponen una factorización a través de \(F\).

\section{Sistemas inversos de supervivientes}

La transición TPK no tiene por qué ser una función. Para tipar correctamente
este hecho, sea \(V_t\) el conjunto finito de estados completos posibles en
la profundidad \(t\), y sea
\[
\mathcal R_t\subseteq V_t\times V_{t+1}
\]
la relación de transición admisible, con todos los campos del estado
enriquecido conservados. Definimos entonces \(X_t\), no como el conjunto de
estados terminales aislados, sino como el conjunto de prefijos admisibles, es
decir, de historias admisibles completas:
\[
X_t=
\Bigl\{
(v_0,\ldots,v_t)\in\prod_{j=0}^{t}V_j:
(v_j,v_{j+1})\in\mathcal R_j\ (0\le j<t),
\ \text{y se satisfacen los cierres hasta }t
\Bigr\}.
\]
Sobre estos objetos la truncación sí es una aplicación canónica:
\[
\begin{aligned}
\rho_{t+1,t}:X_{t+1}&\longrightarrow X_t,\\
(v_0,\ldots,v_t,v_{t+1})&\longmapsto(v_0,\ldots,v_t).
\end{aligned}
\]
Está bien definida precisamente porque todo prefijo de una historia
admisible sigue siendo una historia admisible para los cierres ya impuestos.
Si un filtro utilizase información futura y pudiera invalidar
retroactivamente un prefijo, primero habría que reemplazar \(X_t\) por el
subconjunto de prefijos prolongables; con esa convención, la misma truncación
vuelve a estar bien definida.

Una historia HMT global es una familia de prefijos compatible
\[
x=(x_t)_{t\ge0},
\qquad
x_t\in X_t,
\qquad
\rho_{t+1,t}(x_{t+1})=x_t.
\]
Equivale a una sucesión infinita \((v_0,v_1,\ldots)\) cuyas parejas
consecutivas pertenecen a las relaciones \(\mathcal R_t\) y cuyos prefijos
satisfacen todos los cierres. El espacio de historias es el límite inverso
\[
X_\infty=\varprojlim X_t.
\]
Esta formulación evita imponer una transición local unívoca
\(B_t\mapsto B_{t+1}\). La transición puede ser una relación finita, varias
ramas pueden compartir el mismo estado terminal visible y la unicidad puede
aparecer sólo después de imponer compatibilidad futura. Al conservar el
prefijo completo no se pierde, antes de tiempo, la información que permite
distinguirlas.

\begin{theorem}[Sección coinductiva única]
Supongamos que, para todo \(t\), el conjunto de supervivientes \(S_t\subseteq
X_t\) es no vacío, las restricciones
\(\rho_{t+1,t}:S_{t+1}\to S_t\) son compatibles y cada \(S_t\) contiene un
único prefijo. Entonces \(\varprojlim S_t\) contiene una única historia
global.
\end{theorem}

\begin{proof}
Sea \(x_t\) el único prefijo de \(S_t\). La compatibilidad obliga a
\(\rho_{t+1,t}(x_{t+1})=x_t\), de modo que \((x_t)_t\) pertenece al límite
inverso. Cualquier otra historia tendría en cada nivel un miembro de \(S_t\)
y, por unicidad, coincidiría prefijo a prefijo con \((x_t)_t\).
\end{proof}

El teorema formaliza el paso de la compatibilidad nivel a nivel a una historia
global. No convierte una ventana finita en una conclusión coinductiva: su
aplicación a una familia concreta exige verificar en cada profundidad la no
vacuidad, la compatibilidad de las restricciones y la unicidad de los prefijos
supervivientes.

\section{Conexión nonádica de supervivencia}
\label{sec:conexion-nonadica}

La contracción de cilindros no procede bloque a bloque como una sustitución
determinista. En determinadas fronteras existen varios repartos locales de un
mismo dato agregado; la prolongabilidad decide cuál pertenece a la historia
publicada. La estructura resultante es una conexión sobre las nueve clases de
fase: al completar una vuelta regresa la clase de fase, pero el estado de la
fibra ha cambiado. Esta sección construye la ventana finita completa en la que
se observa ese retorno.

El dominio de este cálculo local es el registro de transporte generado por la
composición APP--TRIT--TPK. Sus estados iniciales, su matriz entera de
propagación y sus relaciones de frontera son publicaciones internas previas de
la transición enriquecida \(\mathcal U_t\), conforme al registro orientado de
doce transiciones y a la reconstrucción única
\(L_i=X_i^{-1}Y_i\). Esta sección los recibe como resultados ya demostrados,
no como premisas externas ni como datos convencionales. Todas sus afirmaciones
son exactas sobre esas salidas. El teorema coinductivo anterior extiende la
misma construcción a profundidad arbitraria y conserva no vacuidad,
compatibilidad de truncamientos y unicidad de los prefijos supervivientes.

\subsection{Prefijo de \(30\) trits y 1.ª frontera}

Escribimos \(u_k^\chi\in\FF_3^6\), con
\(\chi\in\{\pi,e,\varphi\}\), para el bloque visible del canal \(\chi\) a
profundidad \(k\). Los cinco primeros bloques forman tres prefijos de treinta
trits:
\begin{align*}
w_{30}^{\pi}&=010211|012222|010211|002111|110221,\\
w_{30}^{e}&=201101|121221|102011|012222|102011,\\
w_{30}^{\varphi}&=121200|112202|121020|010210|010200.
\end{align*}
El sexto bloque es la primera frontera posterior a esos prefijos:
\[
R_{36}=
\begin{pmatrix}
u_5^\pi\\ u_5^e\\ u_5^\varphi
\end{pmatrix}
=
\begin{pmatrix}
222220\\021222\\102011
\end{pmatrix}.
\]
La notación \(R_{36}\) indica que se han expuesto treinta y seis trits por
canal. No designa un estado terminal: el registro que se reconstruye más abajo
contiene veinte bloques, y la recurrencia admite continuaciones ulteriores
cuando se suministran los estados de Hensel correspondientes.

\subsection{Eje de autoescala y fibra residual}

\ifdefined\HMTAxialTheoremDefined
La descomposición axial de la firma y la identificación del eje
\(\varphi\) frente a la fibra residual \(\pi/e\) ya se demostraron en el
\cref{thm:eje-espejo-nonadico}. Aquí se conserva su consecuencia para la
ventana nonádica sin duplicar el teorema ni la prueba.
\else
Para
\[
B_k=(u_k^\pi,u_k^e,u_k^\varphi)\in(\FF_3^6)^3
\]
definimos, coordenada a coordenada,
\[
q_k=u_k^\pi+u_k^e+u_k^\varphi,
\qquad
a_k=u_k^\pi+u_k^e-u_k^\varphi.
\]

\begin{theorem}[Descomposición axial de la firma de frontera]
\label{thm:eje-espejo-nonadico}
El par \((q_k,a_k)\) determina unívocamente el canal axial y el agregado de
los otros dos canales:
\[
u_k^\varphi=2(q_k-a_k),
\qquad
u_k^\pi+u_k^e=q_k-u_k^\varphi.
\]
En consecuencia, una fibra de la aplicación
\[
q_{\rm ax}(u^\pi,u^e,u^\varphi)
=(u^\pi+u^e,u^\varphi)
\]
sólo puede conservar ambigüedad en el reparto interno entre \(u^\pi\) y
\(u^e\).
\end{theorem}

\begin{proof}
La resta de las dos ecuaciones definitorias da
\(q_k-a_k=2u_k^\varphi\). El inverso de \(2\) en \(\FF_3\) vuelve a ser
\(2\), de donde
\(u_k^\varphi=2(q_k-a_k)\). Sustraer este vector de \(q_k\) produce
\(u_k^\pi+u_k^e\). La operación no recupera cada sumando por separado; por
ello la única fibra residual es el conjunto de repartos del agregado fijo.
\end{proof}

En la terminología geométrica del modelo, \(\varphi\) constituye el eje
fijado por la firma y el par \(\pi/e\) la fibra residual. La expresión no
postula una simetría analítica de las constantes reales: nombra la
descomposición lineal exacta que acaba de demostrarse en \(\FF_3^6\).
\fi

\subsection{Órbita completa de las \(9\) fases consecutivas del Topological Phase Kernel}

Fijamos la fase absoluta mediante
\[
g(k)=
\begin{cases}
k\bmod9,&k\not\equiv0\pmod9,\\
9,&k\equiv0\pmod9.
\end{cases}
\]
Como \(w_{30}\) ocupa \(k=0,\ldots,4\), la primera vuelta completa posterior
al prefijo recorre \(k=5,\ldots,13\), con fases
\(5,6,7,8,9,1,2,3,4\). La enumeración de los \(729=3^6\) candidatos de
firma agregada en cada nivel produce la siguiente ventana:
\begin{center}
\small
\begin{tabular}{@{}ccrrlll@{}}
\toprule
\(k\)&\(g(k)\)&locales&seleccionados&\(u_k^\pi\)&\(u_k^e\)&\(u_k^\varphi\)\\
\midrule
5&5&1&1&222220&021222&102011\\
6&6&1&1&111201&201202&221021\\
7&7&3&1&212121&222210&200221\\
8&8&3&1&200121&212212&110110\\
9&9&2&1&100100&020112&010122\\
10&1&1&1&101222&112221&112002\\
11&2&1&1&022212&110001&100021\\
12&3&1&1&012012&202222&000120\\
13&4&1&1&111210&112102&211122\\
\bottomrule
\end{tabular}
\end{center}
La columna «locales» cuenta los estados de entrada admisibles en el nivel
indicado; no cuenta las salidas de la relación de frontera. Esta distinción
es esencial en \(k=9\): los dos estados locales de entrada preceden a una
fibra de tres salidas del estado seleccionado.

\subsection{Retorno 9→1 de la fase nonádica con cambio de estado y avance de memoria}

Sea
\[
 \pi_B:V_k\longrightarrow\mathcal B:=(\FF_3^6)^3
\]
la proyección al triple visible. La relación de transporte sobre estados
completos induce, una vez fijados los datos de fibra del registro, una relación
visible \(\MonNine^B\). El estado seleccionado que entra en la fase novena es
\[
B_9=(100100\mid020112\mid010122).
\]
Su imagen relacional consta exactamente de
\[
\begin{aligned}
B_{10}^{(0)}&=(101222\mid112221\mid112002),\\
B_{10}^{(1)}&=(102221\mid111222\mid112002),\\
B_{10}^{(2)}&=(111221\mid102222\mid112002).
\end{aligned}
\]
Las tres salidas tienen
\[
q_{10}=022112,
\qquad a_{10}=101111,
\]
y, por el \cref{thm:eje-espejo-nonadico}, comparten
\[
u_{10}^\varphi=112002,
\qquad u_{10}^\pi+u_{10}^e=210110.
\]
Difieren exclusivamente en las tres distribuciones publicadas de ese agregado
entre \(\pi\) y \(e\).

\begin{theorem}[Selección por dos niveles de prolongabilidad]
\label{thm:seleccion-novena-fase}
En la ventana calibrada, las tres salidas anteriores admiten una prolongación
de profundidad uno. Al imponer también el bloque siguiente
\[
B_{11}=(022212\mid110001\mid100021)
\]
y la frontera decimal \((502,662,638)\), sólo
\(B_{10}^{(0)}\) admite una prolongación del mismo estado completo. Por tanto
\[
\#\operatorname{Surv}^{\rm cal}_{1}(B_9)=3,
\qquad
\#\operatorname{Surv}^{\rm cal}_{2}(B_9)=1.
\]
\end{theorem}

\begin{proof}
Las ternas decimales asociadas a las tres salidas son, respectivamente,
\[
(279,352,365),\quad(280,351,365),\quad(282,350,365).
\]
Para una palabra ternaria de longitud \(L\), de entero asociado \(N\), y un
prefijo de \(b\) ternas, de entero asociado \(D\), la compatibilidad de los
dos cilindros equivale a las dos desigualdades enteras
\[
N1000^b<(D+1)3^L,
\qquad
(N+1)1000^b>D3^L.
\]
Las seis desigualdades de cada salida se satisfacen en el nivel actual. Al
anexar \(B_{11}\) y \((502,662,638)\), vuelven a satisfacerse para los tres
canales de \(B_{10}^{(0)}\). Para \(B_{10}^{(1)}\) y \(B_{10}^{(2)}\), se
mantiene la desigualdad del canal \(\varphi\), pero fallan las de \(\pi\) y
\(e\). La evaluación se efectúa con enteros en el certificado finito correspondiente; no
interviene tolerancia numérica. En consecuencia, el conjunto de candidatos
se reduce de tres a uno al pasar de profundidad uno a profundidad dos.
\end{proof}

Las otras dos palabras pueden preceder a una misma cola ternaria sólo si se
cambia la extensión profunda de Hensel. No son, por ello, prolongaciones del
mismo estado completo del registro: el residuo visible coincide en parte, pero
la genealogía de la fibra es distinta.

\subsection{Retorno de fase con cambio de estado}

La fase satisface \(g(k+9)=g(k)\). El estado visible, en cambio, no se
reinicializa. En particular,
\[
B_1=(012222\mid121221\mid112202),
\qquad
B_{10}=(101222\mid112221\mid112002),
\]
de modo que \(g(1)=g(10)=1\) pero \(B_1\ne B_{10}\). La comprobación de los
nueve pares \(k\mapsto k+9\), para \(k=1,\ldots,9\), da además
\[
u_{k+9}^\varphi\ne u_k^\varphi,
\qquad
u_{k+9}^\pi+u_{k+9}^e\ne u_k^\pi+u_k^e.
\]
Por consiguiente, la ventana satisface exactamente
\[
\boxed{g(k+9)=g(k),\qquad B_{k+9}\ne B_k
\quad(1\le k\le9).}
\]
Este es un retorno de fase con holonomía visible no trivial. La denominación
\emph{monodromía nonádica} se refiere a esta acción de retorno sobre la
fibra enriquecida; no introduce una décima fase. Para promover la observación
finita a una monodromía global se debe construir la acción a toda profundidad
y demostrar la compatibilidad coinductiva exigida al comienzo de la sección.

\subsection{La rama de \(20\) bloques}

El retorno anterior no clausura la emisión. Los veinte bloques que el
certificado reconstruye desde los tres estados enteros publicados son
\begin{align*}
\pi:\;&010211|012222|010211|002111|110221|222220|111201|\\
&212121|200121|100100|101222|022212|012012|111210|\\
&121011|200220|120210|000101|022010|020111,\\[2pt]
e:\;&201101|121221|102011|012222|102011|021222|201202|\\
&222210|212212|020112|112221|110001|202222|112102|\\
&102010|022020|222002|012010|001112|022212,\\[2pt]
\varphi:\;&121200|112202|121020|010210|010200|102011|221021|\\
&200221|110110|010122|112002|100021|000120|211122|\\
&202200|202110|221000|100102|111122|020010.
\end{align*}
Así, \(w_{30}\) es el prefijo y \(R_{36}\) la primera frontera de una rama
explícitamente continuada; ninguno de los dos objetos sustituye al estado
enriquecido que transporta esa continuación.

\subsection{Decodificación exacta de \(12\) ternas decimales mediante el lector nonádico}

Sea \(N_{\chi,k}\) el entero representado en base tres por los primeros \(k\)
bloques del canal \(\chi\). Con \(k=13\), la palabra tiene \(78\) trits y
define
\[
D_{\chi,12}
=\left\lfloor\frac{1000^{12}N_{\chi,13}}{3^{78}}\right\rfloor.
\]
La escritura de \(D_{\chi,12}\) en doce bloques de tres cifras es
\[
\begin{array}{c|rrrrrrrrrrrr}
\pi&141&592&653&589&793&238&462&643&383&279&502&884\\
e&718&281&828&459&045&235&360&287&471&352&662&497\\
\varphi&618&033&988&749&894&848&204&586&834&365&638&117.
\end{array}
\]
La identidad es entera: equivale a
\[
D_{\chi,12}3^{78}
\le 1000^{12}N_{\chi,13}
<(D_{\chi,12}+1)3^{78}.
\]
Por tanto, los dos cilindros se intersecan en cada canal y la operación de
lectura produce exactamente las treinta y seis cifras publicadas. Este
resultado es una consecuencia del registro generado. Las construcciones
históricas que emplearon expansiones convencionales pertenecen exclusivamente
a la comparación y al control posterior: no intervienen en la selección
orbital, en la transición \(\mathcal U_t\), en la reconstrucción de
\(L_0,L_1\) ni en el calendario \(0011\). La identidad certifica la lectura
arquimediana de los estados producidos por APP--TRIT--TPK y conserva intacta
su procedencia generativa independiente de valores objetivo.

\subsection{Calendario de reapertura de cilindros y periodicidad de la prolongación compatible}

La conversión de un bloque de seis trits a ternas decimales tiene pendiente
media
\[
\lambda=6\log_{1000}3\in(0,1),
\qquad k(t)=\lfloor\lambda(t+1)\rfloor.
\]
Un índice \(t\) es una reapertura sin nueva terna cuando
\(k(t+1)=k(t)\). Sea \(\delta=1-\lambda\).

\begin{theorem}[Calendario de Beatty de la lectura ternario--decimal]
\label{thm:calendario-beatty}
El conjunto de reaperturas es
\[
\mathcal T=
\left\{\left\lceil\frac{m}{\delta}\right\rceil-2:m\ge1\right\}.
\]
Sus primeros elementos son
\[
20,42,64,86,108,130,151,173,195,217,\ldots,
\]
y dos reaperturas consecutivas distan \(21\) o \(22\) bloques.
\end{theorem}

\begin{proof}
La cantidad \(\lambda=\log_{1000}729\) es irracional: una igualdad
\(\lambda=p/q\) implicaría \(729^q=1000^p\), imposible por factorización
prima. También \(\delta\) es irracional, por lo que
\[
k(t)=(t+1)-\lceil\delta(t+1)\rceil.
\]
La igualdad \(k(t+1)=k(t)\) ocurre exactamente cuando la función techo del
segundo término aumenta en una unidad. La \(m\)-ésima ocasión se obtiene
resolviendo
\(\delta(t+1)<m<\delta(t+2)\), lo cual da
\(t=\lceil m/\delta\rceil-2\). Las diferencias de una sucesión de Beatty
pertenecen a
\[
\left\{\left\lfloor\delta^{-1}\right\rfloor,
\left\lceil\delta^{-1}\right\rceil\right\}=\{21,22\}.\qedhere
\]
\end{proof}

\subsection{Recurrencia Bockstein--Hensel de residuo y cociente}

La memoria transportada por la rama procede de la sucesión exacta
\[
0\longrightarrow\ZZ^6\xrightarrow{\,\times3\,}\ZZ^6
\longrightarrow\FF_3^6\longrightarrow0
\]
y de la matriz unimodular
\[
A=
\begin{pmatrix}
2&2&2&1&2&1\\
2&1&2&2&1&1\\
1&0&0&1&0&2\\
0&0&1&1&1&0\\
2&2&2&0&2&1\\
2&1&2&1&0&0
\end{pmatrix},
\qquad\det A=1.
\]
Para un estado fila \(x_k^\chi\in\ZZ^6\), se define
\[
y_k^\chi=x_k^\chi A,
\qquad
u_k^\chi=y_k^\chi\bmod3\in\{0,1,2\}^6,
\qquad
x_{k+1}^\chi=\frac{y_k^\chi-u_k^\chi}{3}.
\]

\begin{proposition}[Conservación exacta de la memoria de cociente]
\label{prop:memoria-bockstein-hensel}
Fijados \(x_k^\chi\) y \(A\), el par
\((u_k^\chi,x_{k+1}^\chi)\) es único y satisface
\[
x_k^\chi A=u_k^\chi+3x_{k+1}^\chi.
\]
Recíprocamente, como \(A\in\operatorname{GL}_6(\ZZ)\), el par determina
unívocamente \(x_k^\chi\).
\end{proposition}

\begin{proof}
La división euclídea de cada coordenada de \(x_k^\chi A\) por tres produce
un residuo único en \(\{0,1,2\}\) y un cociente entero único. Esto prueba
la primera afirmación y la identidad. Puesto que \(\det A=1\), la inversa
\(A^{-1}\) tiene entradas enteras; por tanto
\[
x_k^\chi=(u_k^\chi+3x_{k+1}^\chi)A^{-1}
\]
recupera el estado anterior sin ambigüedad.
\end{proof}

El residuo \(u_k^\chi\) es la parte visible y el cociente
\(x_{k+1}^\chi\) la memoria 3-ádica transportada. El certificado reproduce
sesenta pasos coordenados ---veinte por canal---, verifica \(AA^{-1}=I\) con
aritmética entera y coteja todos los bloques con el registro. La denominación
Bockstein--Hensel alude aquí a esta separación residuo--cociente y a su
iteración 3-ádica; no se identifica sin datos adicionales con un operador de
Bockstein cohomológico.

\section{Del prefijo a un punto real}

Asociemos a cada prefijo compatible \(x_t\in X_t\) un intervalo cerrado
\(I_t(x_t)\subset\mathbb R\). La interpretación geométrica es la de un
cilindro: todos los futuros que comparten el mismo prefijo se proyectan dentro
de \(I_t\).

\begin{theorem}[Convergencia arquimediana de intervalos encajados]
Sea \((x_t)_{t\ge0}\) una historia compatible. Si
\[
I_{t+1}(x_{t+1})\subseteq I_t(x_t)
\]
y
\[
\operatorname{diam} I_t(x_t)\longrightarrow0,
\]
entonces existe un único \(r\in\mathbb R\) tal que
\[
\bigcap_{t\ge0}I_t(x_t)=\{r\}.
\]
\end{theorem}

\begin{proof}
Es el principio de intervalos encajados. Los extremos izquierdos forman una
sucesión creciente acotada y los derechos una decreciente acotada. Sus límites
existen; la diferencia entre ellos está acotada por el diámetro, que tiende a
cero. Ambos límites coinciden en un único punto contenido en todos los
intervalos.
\end{proof}

Sea \(X_\infty^{\rm Arch}\subseteq X_\infty\) el dominio de historias para
las que los cilindros están encajados y sus diámetros tienden a cero. Si esas
dos propiedades forman parte de la admisibilidad global, entonces
\(X_\infty^{\rm Arch}=X_\infty\); si no, esta restricción de dominio es
necesaria. Definimos la evaluación
\[
\operatorname{ev}_{\mathbb R}:X_\infty^{\rm Arch}\longrightarrow\mathbb R,
\qquad
\operatorname{ev}_{\mathbb R}(x)=
\text{el único punto de }\bigcap_{t\ge0}I_t(x_t).
\]
La unicidad del valor no obliga a la unicidad de la historia. Dos genealogías
\(x\ne y\) pueden satisfacer
\[
\operatorname{ev}_{\mathbb R}(x)
=\operatorname{ev}_{\mathbb R}(y).
\]
Definimos la equivalencia de lectura
\[
x\sim_{\rm ev}y
\quad\Longleftrightarrow\quad
\operatorname{ev}_{\mathbb R}(x)
=\operatorname{ev}_{\mathbb R}(y).
\]

\begin{theorem}[Cociente de la evaluación arquimediana]
La aplicación
\[
\begin{aligned}
\overline{\operatorname{ev}}_{\mathbb R}:
X_\infty^{\rm Arch}/\!\sim_{\rm ev}
&\longrightarrow
\operatorname{Im}(\operatorname{ev}_{\mathbb R})\subseteq\mathbb R,\\
[x]&\longmapsto\operatorname{ev}_{\mathbb R}(x)
\end{aligned}
\]
es una biyección. En particular, el cociente queda identificado con todo
\(\mathbb R\) si y sólo si la evaluación
\(\operatorname{ev}_{\mathbb R}\) es sobreyectiva.
\end{theorem}

\begin{proof}
La definición de \(\sim_{\rm ev}\) hace que el valor sea independiente del
representante, así que la aplicación está bien definida. Es inyectiva porque
dos clases con el mismo valor son, por definición, una sola clase; es
sobreyectiva sobre la imagen por la definición misma de imagen. La última
afirmación equivale a
\(\operatorname{Im}(\operatorname{ev}_{\mathbb R})=\mathbb R\).
\end{proof}

El teorema de cobertura del lector arquimediano construye precisamente esos
prefijos: para todo intervalo compacto \(K\subset\mathbb R\),
\(\operatorname{ev}_{\mathbb R}(\Omega_{\infty,K})=K\), y los intervalos
\([-m,m]\) forman una familia dirigida cuya unión es \(\mathbb R\). Por
consiguiente, cada real posee al menos una historia HMT compatible y la
afirmación demostrada es
\[
X_\infty^{\rm Arch}/\!\sim_{\rm ev}
\cong\operatorname{Im}(\operatorname{ev}_{\mathbb R})
=\mathbb R.
\]
El valor arquimediano es, por tanto, el cociente que identifica las fibras de
la evaluación; el número HMT enriquecido conserva además la clase de
trayectoria, el acarreo y la orientación. La identificación no se sigue de la
mera formación abstracta del cociente: se sigue del teorema de cobertura ya
demostrado en la construcción discreta del continuo.

El descenso algebraico es otro problema y debe resolverse operación por
operación. Supongamos que se han construido dos operaciones candidatas
\(\oplus\) y \(\odot\) sobre \(X_\infty^{\rm Arch}\), cerradas en ese
dominio. La suma desciende al cociente si y sólo si
\[
x\sim_{\rm ev}x',\ y\sim_{\rm ev}y'
\Longrightarrow
x\oplus y\sim_{\rm ev}x'\oplus y',
\]
mientras que el producto desciende, de manera independiente, si y sólo si
\[
x\sim_{\rm ev}x',\ y\sim_{\rm ev}y'
\Longrightarrow
x\odot y\sim_{\rm ev}x'\odot y'.
\]
La primera congruencia no implica la segunda. Además, para que las operaciones
inducidas sean precisamente la suma y el producto ordinarios del cociente
arquimediano se necesita probar por separado
\[
\operatorname{ev}_{\mathbb R}(x\oplus y)
=\operatorname{ev}_{\mathbb R}(x)+\operatorname{ev}_{\mathbb R}(y),
\]
\[
\operatorname{ev}_{\mathbb R}(x\odot y)
=\operatorname{ev}_{\mathbb R}(x)\operatorname{ev}_{\mathbb R}(y),
\]
y que la imagen de la evaluación sea cerrada bajo la operación correspondiente.
Sólo después de estos certificados el cociente hereda la estructura de
subanillo ---o de subcuerpo, si se construyen también opuestos e inversos---
de \(\mathbb R\). La biyección conjuntista del teorema anterior es formal;
el contenido específico de HMT reside en construir la evaluación, caracterizar sus
fibras, demostrar cobertura y establecer cada ley de descenso.

\section{Representación de Weyl--Heisenberg asociada al toro nonádico}

La forma de área orientada del toro nonádico es
\[
\sigma((a,b),(c,d))=ad-bc\pmod9.
\]
Esta forma alternante es la que se obtiene, con la orientación elegida, al
integrar la curvatura mixta de APP sobre paralelogramos del toro. Sea
\(\omega=\exp(2\pi i/9)\) y
\(\mathcal H_9=\mathbb C^9\), con base \(\{|n\rangle:n\in\mathbb Z_9\}\).
Definimos
\[
X|n\rangle=|n+1\rangle,
\qquad
Z|n\rangle=\omega^n|n\rangle.
\]
Entonces
\[
ZX=\omega XZ,
\]
y, más generalmente,
\[
(X^aZ^b)(X^cZ^d)
=\omega^{bc}X^{a+c}Z^{b+d}.
\]
Por tanto
\[
 c\bigl((a,b),(c,d)\bigr)=bc\pmod9
\]
es el \(2\)-cociclo polarizado que aparece en esta sección de la extensión
central. Definimos explícitamente
\[
 \operatorname{Heis}_9
 =\{\omega^kX^aZ^b:a,b,k\in\ZZ/9\ZZ\}.
\]
Es un grupo de orden \(9^3=729\). Se trata de una extensión central asociada
al grupo de traslaciones del soporte \(X_9\) de APP, no de un subgrupo del
grafo ni de la categoría de caminos de APP\@.
El conmutador de los dos transportes es
\[
(X^aZ^b)(X^cZ^d)(X^aZ^b)^{-1}(X^cZ^d)^{-1}
=\omega^{bc-ad}I.
\]
En consecuencia, su exponente es \(-\sigma((a,b),(c,d))\). El cociclo
polarizado \(c\) y la forma \(\sigma\) no son el mismo objeto: la segunda es,
con este convenio, la opuesta de la antisimetrización
\(c(u,v)-c(v,u)\). La fase de conmutador es así la inversa de la fase de
Wilson para la orientación fijada.

Para construir observables, tomamos
\[
Q_{\rm H}=\operatorname{diag}(0,1,\ldots,8),
\qquad
F_{jk}=\frac13\omega^{jk},
\qquad
P_{\rm H}=F^\dagger Q_{\rm H}F.
\]
Los operadores \(Q_{\rm H},P_{\rm H}\) son autoadjuntos. Sus bases propias son mutuamente
insesgadas porque \(|F_{jk}|^2=1/9\). Para todo vector unitario
\(\psi\in\mathcal H_9\), la desigualdad de Robertson da
\[
\operatorname{Var}_\psi(Q_{\rm H})\,
\operatorname{Var}_\psi(P_{\rm H})
\ge
\frac14\left|\langle\psi,[Q_{\rm H},P_{\rm H}]\psi\rangle\right|^2,
\]
y la cota entrópica para bases mutuamente insesgadas produce
\[
H_{Q_{\rm H}}(\psi)+H_{P_{\rm H}}(\psi)\ge\log9.
\]
Aquí \(H\) es la entropía de Shannon con logaritmo natural de las
distribuciones espectrales correspondientes. Las relaciones construidas se
resumen, sin atribuirles una causalidad física adicional, por
\[
\begin{gathered}
\text{forma alternante }(X_9,\sigma)
\longrightarrow
\operatorname{Heis}_9,\\
\operatorname{Heis}_9\longrightarrow(X,Z)
\longrightarrow(Q_{\rm H},P_{\rm H})
\longrightarrow\text{incertidumbre finita}.
\end{gathered}
\]
Éste es un modelo algebraico finito de operadores; por sí solo no especifica
una dinámica física ni unidades de medida.

\section{Dinámica cíclica sobre \texorpdfstring{\(\mathbb Z_9\times\mathbb Z_{12}\)}{Z9 x Z12} y cociente predictivo mínimo}

El sistema dinámico finito \(\mathfrak C_{9,12}\) permite determinar la memoria
mínima necesaria para que la evolución observada sea determinista. Sea
\[
X=\{(b,m):b\in\mathbb Z/9\mathbb Z, 0\le m<12\}
\]
con sucesor
\[
S(b,m)=
\begin{cases}
(b,m+1),&m<11,\\
(b+1\bmod9,0),&m=11.
\end{cases}
\]
Las etiquetas marginales son
\[
\beta(b,m)=4b\pmod{12},
\qquad
d(b,m)=1+(m+\beta(b,m)\pmod{12}).
\]
Para una aplicación de observación \(q\), la palabra futura inclusiva de horizonte \(h\) es
\[
Q_h^q(x)=\bigl(q(x),q(Sx),\ldots,q(S^hx)\bigr).
\]

\begin{theorem}[Cociente predictivo mínimo]
Sean \(X\) finito, \(T:X\to X\) y \(q:X\to V\). Definamos
\[
\operatorname{Eq}(f):=\{(x,y)\in X^2:f(x)=f(y)\},
\qquad
E_h=\operatorname{Eq}(Q_h^q),
\qquad
E_\infty=\bigcap_{k\ge0}\operatorname{Eq}(q\circ T^k).
\]
Entonces
\(E_{h+1}=E_h\cap(T\times T)^{-1}(E_h)\), la sucesión estabiliza en tiempo
finito y \(X/E_\infty\) es el cociente con menor número de estados que
refina a \(q\) y admite una dinámica determinista inducida.
\end{theorem}

\begin{proof}
\(E_h\) exige igualdad de observaciones en los tiempos \(0,\ldots,h\), mientras
\((T\times T)^{-1}(E_h)\) exige igualdad en \(1,\ldots,h+1\). Su intersección es
\(E_{h+1}\). Como \(X\) es finito, el número de clases no puede crecer
indefinidamente. En el primer nivel estable, \(T\) respeta las clases y
desciende. Toda equivalencia \(T\)-compatible contenida en
\(\operatorname{Eq}(q)\)
identifica sólo estados con el mismo futuro completo, por lo que está
contenida en \(E_\infty\); de ahí la propiedad minimal.
\end{proof}

La evaluación exhaustiva de \(\mathfrak C_{9,12}\) da:
\begin{center}
\small
\begin{tabular}{@{}lrrrr@{}}
\toprule
observable & valores instantáneos & índice \(h\) & observaciones & cociente estable\\
\midrule
\(\beta\) & 3 & 11 & 12 & \(C_{36}\), fibra 3\\
\(d\) & 12 & 8 & 9 & \(C_{36}\), fibra 3\\
\((\beta,d)\) & 36 & 0 & 1 & \(C_{36}\), fibra 3\\
\bottomrule
\end{tabular}
\end{center}
En los dos observables marginales,
\[
\operatorname{Eq}(Q_{11}^{\beta})
=\operatorname{Eq}(Q_8^d)
=\operatorname{Eq}(\pi_{36}),
\qquad
\pi_{36}(b,m)=(b\bmod3,m).
\]
El factor \(36\) no se introduce como objetivo. Aparece como el cociente
predictivo mínimo de cada marginal. La fase \(m\) es una coordenada de estado
necesaria para predecir esos futuros; esta afirmación matemática no presupone
una interpretación causal física. La transformación \(b\mapsto b+3\) es una
simetría invisible para esos observables en todo tiempo.

Como sistema dinámico abstracto, el cociente es un ciclo de treinta y seis
estados. Escribir \(C_{36}\) no elige, sin embargo, un origen distinguido:
para obtener una coordenada \(\mathbb Z/36\mathbb Z\) hay que fijarlo. Tampoco
convierte automáticamente esos estados en los treinta y seis pares de una
nónada. En efecto, \(\operatorname{Aut}J(9,2)\cong S_9\). Si los dos puntos de
un par están en ciclos distintos de longitudes \(r,s\) de una permutación, la
órbita del par tiene longitud \(\operatorname{mcm}(r,s)\), con \(r+s\leq9\),
y por tanto longitud a lo sumo \(20\). Si están en un mismo ciclo de longitud
\(r\leq9\), la órbita tiene longitud \(r\), salvo el caso antipodal, en que
tiene longitud \(r/2\). Ningún automorfismo induce, pues, una órbita de longitud
\(36\) sobre los pares. Por tanto, el paso
\[
C_{36}\longleftrightarrow \binom{[9]}2
\]
permanece como identificación calibrada, no como identificación equivariante natural.
Incluso una ordenación hamiltoniana de los pares probaría conservación local
de incidencia paso a paso, no un automorfismo global ni la canonicidad del
diccionario.

\section{Número HMT como historia compatible en \(X_\infty\): identidad genealógica, estructura de lectura y evaluación arquimediana}

Los resultados precedentes determinan la siguiente instancia de la definición
rectora. En la notación de esta sección, un número HMT es una historia
compatible
\[
x\in X_\infty=\varprojlim X_t.
\]
Un lector \(L\) es estructura adicional: el par \((x,L)\) constituye una
presentación leída o una medición, pero no altera el tipo de identidad de
\(x\). Cuando \(x\) pertenece a \(X_\infty^{\rm Arch}\), su valor
arquimediano es la intersección puntual de los cilindros; el conjunto de
valores así construido es
\(\operatorname{Im}(\operatorname{ev}_{\mathbb R})\subseteq\mathbb R\), y
su identificación con toda la recta exige el teorema adicional de cobertura.
Su memoria mínima relativa a un observable es el cociente predictivo de
futuros; y la forma alternante del soporte suma--producto admite la
representación de Weyl--Heisenberg finita construida arriba. El valor, la clase de trayectoria y la incertidumbre quedan
definidos por tres operaciones sobre dominios especificados de una misma
arquitectura de transporte.
